\section{Introduction}
\label{sec:introduction}

Teaching the transient and steady-state behavior of analog circuits is difficult precisely
because the phenomena involved -- charge accumulating on a capacitor plate, flux linking an
inductor's turns, a memristor's internal state drifting with the integral of the current
through it -- are invisible. Students are asked to reason about quantities they cannot see,
using differential equations whose solutions they are often meeting for the first time in
the same course. The traditional response is to pair the mathematics with a laboratory
bench and, increasingly, with software simulators such as SPICE~\cite{nagel1973spice},
LTspice, or browser-based tools such as the Falstad circuit simulator and PhET's Circuit
Construction Kit. These tools are effective, but they present the circuit as a schematic on
a two-dimensional canvas -- an abstraction one level removed from the physical act of
connecting components that a real breadboard or benchtop still preserves.

Sandbox games occupy an unusual middle ground. Minecraft in particular has an established
track record as an informal STEM learning environment, precisely because building something
inside it -- a redstone logic circuit, a working calculator, a piston-driven
elevator -- requires the same kind of systems thinking as an engineering task, while
remaining playful and low-stakes~\cite{nebel2016mining}. Minecraft's own redstone system is
already a widely used entry point into digital logic. There is, however, no equivalent for
\emph{analog} circuit theory inside the base game: redstone signals are effectively boolean,
propagated instantaneously (subject to a fixed per-block delay), with no notion of voltage,
current, capacitance, or continuous time.

This paper describes Mine Memristors, a modification (``mod'') for the Fabric mod loader
that adds a full set of analog circuit components to Minecraft: resistors, capacitors,
inductors, memristors, DC and time-varying voltage sources, an ammeter, wire, an explicit
ground reference block, and a handheld multi-channel oscilloscope probe. Placing and wiring
these blocks together builds a real circuit topology in the game world; a modified nodal
analysis solver, structurally the same technique used by SPICE-family
simulators~\cite{ho1975mna}, is re-solved every game tick (20~Hz) to drive every block's
displayed voltage, current, and internal state. Nothing in the mod is a scripted
approximation: an RC charging curve looks like an RC charging curve because the underlying
linear system is actually being solved, and a memristor's resistance drifts according to the
same charge-controlled ordinary differential equation used in the original HP linear-drift
model~\cite{strukov2008missing}, not a lookup table tuned to look plausible.

The remainder of this paper is organized as follows. Section~\ref{sec:related-work}
situates Mine Memristors relative to existing circuit simulators and the literature on
game-based engineering education. Section~\ref{sec:architecture} describes the mod's
software architecture, in particular the topology-construction algorithm used to translate
a player's in-world wiring into a circuit graph, and a subtle correctness bug in an earlier
version of that algorithm that is instructive in its own right.
Section~\ref{sec:circuit-solver} describes the solver itself, including the explicit,
player-placeable ground reference that lets the notion of ``ground'' be taught rather than
assumed. Section~\ref{sec:pedagogy} discusses the pedagogical rationale for each design
choice. Section~\ref{sec:validation} describes how the solver and the surrounding game logic
were verified during development. Section~\ref{sec:limitations} discusses current
limitations and planned classroom evaluation, and Section~\ref{sec:conclusion} concludes.
